% Vietnamese translation for OI Public Library
% Source-File: IOI中国国家候选队论文/2005/王俊/王俊.doc
\documentclass[11pt]{article}
\usepackage{oipl-vn}

\newcommand{\Oh}{\mathcal{O}}

\title{Bàn sơ về ứng dụng của ghép cặp đồ thị hai phía trong thi tin học}
\author{Wang Jun\\Trường Trung học Changjun, Trường Sa, Hồ Nam}
\date{2005}

\begin{document}
\maketitle

\origtitle{Bàn sơ về ứng dụng của ghép cặp đồ thị hai phía trong thi tin học}
\sourcepath{IOI China candidate team papers / 2005 / Wang Jun.doc}

\renewcommand{\abstractname}{Tóm tắt}
\begin{abstract}
Bài viết thông qua phân tích một số bài thi tin học để minh họa các ứng dụng
của ghép cặp đồ thị hai phía. Thông thường, ta phân tích tính chất của một bài
toán tối ưu, xây dựng đồ thị hai phía, rồi giải bài toán gốc bằng cách tìm
ghép cực đại, ghép tốt nhất hoặc các dạng ghép khác trên đồ thị ấy.
\end{abstract}

\noindent\textbf{Từ khóa:} ghép cặp; đồ thị hai phía; trọng số nhỏ nhất; trọng
số lớn nhất; tối ưu hóa.

\section{Dẫn nhập}

Ghép cặp đồ thị hai phía là một lớp thuật toán đồ thị kinh điển trong thi tin
học, và xuất hiện rất rộng trong những năm gần đây. Nếu có thể chia các đối
tượng của đề bài thành hai tập bổ sung cho nhau, đồng thời cần tìm một quan
hệ một--một thỏa điều kiện nào đó giữa hai tập, ta thường có thể trừu tượng
hóa đối tượng và quan hệ thành một đồ thị hai phía, rồi dùng thuật toán ghép
cặp để giải.

Những bài như vậy thường kiểm tra nhiều năng lực: mô hình hóa đề gốc, xây
dựng đồ thị hai phía, thiết kế thuật toán ghép, và tối ưu thuật toán theo tính
chất riêng của bài. Hai ví dụ dưới đây minh họa một số cách dùng ghép cặp
trong thi tin học.

\section{Railway Communication}

\subsection{Mô tả bài toán}

Một quốc gia có $n$ thị trấn và $m$ đường sắt một chiều. Mỗi đường nối hai thị
trấn khác nhau, và hệ thống đường sắt không có chu trình. Cần xác định một số
tuyến chạy tàu thỏa mãn:
\begin{itemize}
  \item mỗi đường sắt thuộc nhiều nhất một tuyến;
  \item mỗi thị trấn được nhiều nhất một tuyến đi qua, kể cả làm điểm đầu hoặc
  điểm cuối;
  \item mỗi thị trấn được ít nhất một tuyến đi qua;
  \item số tuyến càng ít càng tốt;
  \item trong các phương án có số tuyến ít nhất, tổng chi phí đường sắt được
  dùng càng nhỏ càng tốt.
\end{itemize}

\subsection{Phân tích}

Bài toán vừa yêu cầu số tuyến ít nhất, vừa yêu cầu tổng độ dài hoặc chi phí
nhỏ nhất trong điều kiện đó. Ta tách thành hai bước: trước hết chỉ xét bài con
số tuyến ít nhất. Mô hình toán học là: cho một DAG
\[
  G_0=(N_0,A_0),
\]
dùng ít đường đi đơn rời nhau nhất để phủ toàn bộ tập đỉnh $N_0$.

Ta xây dựng đồ thị hai phía $G=(N,A)$. Tạo hai tập đỉnh bổ sung $X$ và $Y$.
Mỗi đỉnh $i\in N_0$ được tách thành một đỉnh $i$ ở $X$ và một đỉnh $i'$ ở
$Y$. Với mỗi cung $(i,j)$ của $G_0$, thêm cạnh $(i,j')$ vào $G$. Nếu trong
$G_0$ chọn cung $(i,j)$ làm một cạnh của đường phủ, thì trong $G$ chọn cạnh
$(i,j')$.

Xét một đỉnh $i$ của $G_0$:
\begin{itemize}
  \item Nếu $i$ là điểm đầu của một đường phủ, nó không có tiền nhiệm; trong
  đồ thị hai phía, không cạnh nào kề $i'$ được chọn.
  \item Nếu $i$ nằm bên trong một đường phủ, nó có một tiền nhiệm và một hậu
  nhiệm; trong đồ thị hai phía, mỗi phía $i$ và $i'$ đều có đúng một cạnh được
  chọn.
  \item Nếu $i$ là điểm cuối, nó không có hậu nhiệm; không cạnh nào kề $i$ được
  chọn.
\end{itemize}

Như vậy bài toán trở thành chọn một số cạnh trong đồ thị hai phía sao cho mỗi
đỉnh kề nhiều nhất một cạnh được chọn. Đây chính là ghép cặp. Số đường phủ
bằng số đỉnh không có hậu nhiệm trong cách phủ, nên muốn số đường phủ nhỏ nhất
thì phải chọn càng nhiều cạnh ghép càng tốt. Do đó bài con là bài ghép cực đại
trong đồ thị hai phía, có thể dùng thuật toán Hungary kinh điển.

Bây giờ xét yêu cầu tổng chi phí nhỏ nhất. Nếu cung $(i,j)$ trong $G_0$ có
trọng số $w_{ij}$, ta gán cùng trọng số cho cạnh $(i,j')$ trong $G$. Bài toán
trở thành: trong số các ghép có kích thước cực đại, tìm ghép có tổng trọng số
nhỏ nhất, tức là ghép cực đại trọng số nhỏ nhất.

Thuật toán KM kinh điển thường giải ghép cực đại trọng số lớn nhất. Ta chuyển
trọng số bằng cách chọn một hằng số lớn $W$ và đặt trọng số mới
\[
  W'_{ij}=W-w_{ij}
\]
cho các cạnh tồn tại; các cạnh không tồn tại được xử lý sao cho không thể được
chọn hoặc có trọng số rất bất lợi. Khi đó ghép trọng số lớn nhất theo $W'$ sẽ
tương ứng với ghép trọng số nhỏ nhất theo $w$, trong khi kích thước ghép vẫn
cực đại. Như vậy bài toán được giải bằng KM.

\subsection{Tiểu kết}

Bài này là một mở rộng của bài phủ đường đi nhỏ nhất trên DAG. Điều kiện quan
trọng là mỗi đỉnh trên một đường phủ có nhiều nhất một tiền nhiệm và một hậu
nhiệm; điều đó tương ứng trực tiếp với tính chất mỗi đỉnh trong ghép cặp chỉ
được ghép với nhiều nhất một đỉnh. Quan hệ một--một là tính chất cốt lõi của
ghép cặp.

\section{Roads}

\subsection{Mô tả bài toán}

Một vương quốc có $n$ thành phố và $m$ con đường. Trong đó có $n-1$ đường đá
và $m-n+1$ đường đất. Giữa hai thành phố bất kỳ tồn tại đúng một đường đi chỉ
gồm đường đá, nên các đường đá tạo thành một cây. Các đường được đánh số duy
nhất: đường đá là $1..n-1$, đường đất là $n..m$. Đường thứ $i$ có chi phí bảo
trì thật là $C_i$.

Nhà vua muốn chỉ bảo trì một số đường sao cho mọi thành phố vẫn liên thông,
và giao cho đại thần chọn tập đường có tổng chi phí nhỏ nhất. Dù vua không
phân biệt đường đá với đường đất, dân chúng thích đường đá hơn. Vì vậy ta muốn
giả mạo chi phí bảo trì thành $D_i$ sao cho cây gồm các đường đá được chọn là
một cây khung nhỏ nhất. Để giảm khả năng bị phát hiện, cần cực tiểu hóa
\[
  \sum_{i=1}^{m}|C_i-D_i|.
\]
Nếu có nhiều cây khung cùng chi phí nhỏ nhất, đại thần có thể chọn cây đường
đá, nên chỉ cần khiến cây đá là một trong các cây khung nhỏ nhất.

\subsection{Điều kiện để cây đá là MST}

Gọi đồ thị gốc là $G_0=(N_0,A_0)$, và gọi cây gồm các cạnh đá là $T$. Với một
cạnh không thuộc cây $u$, thêm $u$ vào $T$ tạo ra đúng một chu trình. Ký hiệu
$P[u]$ là tập các cạnh trên đường đi trong $T$ nối hai đầu mút của $u$.

Theo tính chất cây khung nhỏ nhất, để $T$ là một MST, với mọi cạnh không thuộc
cây $u$ và mọi cạnh cây $v\in P[u]$, cần có
\[
  D_v\le D_u.
\]
Nếu không, thay $v$ bằng $u$ sẽ tạo một cây khung có trọng số nhỏ hơn.

Bất đẳng thức trên luôn có $v$ là cạnh cây và $u$ là cạnh ngoài cây. Trực giác
cho thấy để giảm tổng thay đổi, ta chỉ nên giảm trọng số cạnh cây và tăng
trọng số cạnh ngoài cây. Đặt lượng sửa đổi không âm:
\[
  x_v=C_v-D_v\quad (v\in T),\qquad
  y_u=D_u-C_u\quad (u\notin T).
\]
Khi đó điều kiện $D_v\le D_u$ tương đương
\[
  x_v+y_u\ge C_v-C_u.
\]
Mục tiêu là cực tiểu hóa
\[
  \sum_{v\in T}x_v+\sum_{u\notin T}y_u
\]
dưới các ràng buộc trên.

Dạng bất đẳng thức này rất giống điều kiện nhãn khả thi trong thuật toán KM.
Trong KM, với mỗi cạnh trọng số $w_{ij}$, nhãn hai phía $l_i,r_j$ luôn thỏa
\[
  l_i+r_j\ge w_{ij}.
\]
Vì vậy ta có thể xây dựng một đồ thị hai phía có trọng số để đưa bài toán về
ghép cặp tốt nhất.

\subsection{Mô hình ghép cặp trọng số}

Tạo hai tập đỉnh $X,Y$. Mỗi cạnh cây $v$ được đưa vào phía $X$, mỗi cạnh
ngoài cây $u$ được đưa vào phía $Y$. Nếu $v\in P[u]$, thêm cạnh giữa $v$ và
$u$ với trọng số
\[
  W_{vu}=C_v-C_u.
\]
Các ràng buộc $x_v+y_u\ge W_{vu}$ trở thành đúng dạng nhãn khả thi.

Bản gốc xây dựng thêm các đỉnh và cạnh ảo để biến đồ thị thành đồ thị hai phía
đầy đủ, phù hợp với phiên bản KM dùng ghép hoàn hảo. Khi chạy KM, tổng nhãn
nhỏ nhất tương ứng với tổng sửa đổi nhỏ nhất, và từ nhãn thu được ta khôi phục
được một phương án $D$ hợp lệ. Nói cách khác, bài toán inverse MST này được
chuyển thành bài toán ghép cặp trọng số lớn nhất, qua đối ngẫu của KM.

\subsection{Nếu chỉ cần giá trị tối ưu}

Nếu đề chỉ yêu cầu lượng sửa đổi nhỏ nhất, không cần in phương án $D_i$, ta có
thể khai thác thêm và tránh một số đỉnh ảo.

\paragraph{Cách 1: KM trực tiếp.}
Dùng mô hình ở trên và chạy KM để lấy giá trị ghép lớn nhất. Đây là cách đơn
giản, nhưng vẫn giữ chi phí của đồ thị đầy đủ sau khi thêm đỉnh ảo.

\paragraph{Cách 2: mạng dòng chi phí lớn nhất.}
Vì các đỉnh và cạnh ảo chủ yếu phục vụ điều kiện ghép hoàn hảo của KM, nếu
chỉ cần giá trị, ta có thể bỏ chúng và xây dựng mạng:
nguồn nối tới các cạnh cây, cạnh cây nối tới cạnh ngoài cây khi quan hệ
$v\in P[u]$ tồn tại, cạnh ngoài cây nối tới đích. Cạnh giữa hai phía mang chi
phí tương ứng $W_{vu}$ và dung lượng $1$. Khi đó ghép trọng số lớn nhất tương
ứng với dòng cực đại chi phí lớn nhất. Dùng Bellman--Ford tìm đường tăng chi
phí lớn nhất cho ta một thuật toán đúng nhưng còn chậm.

\paragraph{Cách 3: dùng SPFA trong mạng dư.}
Nút thắt của cách 2 là tìm đường tăng chi phí lớn nhất trong mạng dư. Có thể
dùng SPFA, tức phiên bản hàng đợi của Bellman--Ford, để tăng tốc trong thực
tế. SPFA lưu giá trị đường đi tốt nhất tạm thời, đưa các đỉnh cần tối ưu vào
hàng đợi, và lặp cho tới khi hàng đợi rỗng. Khi không có chu trình âm hoặc
dương bất lợi theo dạng bài, thuật toán hội tụ. Trong các dữ liệu thông thường,
số lần một đỉnh vào hàng đợi thường nhỏ, nên tốc độ thực tế tốt hơn nhiều.

\paragraph{Cách 4: chuyển trọng số từ cạnh sang đỉnh.}
Quay lại bản chất ghép cặp, bản gốc nêu một cách nhìn khác: thay vì đặt trọng
số trên cạnh, chuyển trọng số sang đỉnh phía $Y$. Tạo đồ thị hai phía $G'$,
trong đó các đỉnh tương ứng với cạnh cây và cạnh ngoài cây; thêm cạnh đồng
nhất $(i,i)$ và các cạnh thể hiện quan hệ chu trình. Trọng số gắn với đỉnh
$Y$. Khi đó mục tiêu trở thành chọn các đỉnh $Y$ có trọng số nhỏ để được ghép.

Thuật toán sắp các đỉnh phía $Y$ theo trọng số không giảm, rồi lần lượt thử
tìm ghép cho từng đỉnh. Duy trì tập các đỉnh phía $X$ còn có thể ghép và các
đường tăng tương ứng; mỗi khi tìm được một ghép mới thì cập nhật tập khả ghép.
Nhờ chỉ làm việc trên phần đồ thị liên quan tới các đường tăng hiện tại, thuật
toán có thể tốt hơn các cách trực tiếp. Bản gốc phân tích rằng độ phức tạp
được cải thiện đáng kể so với KM đơn giản, dù còn có thuật toán tốt hơn trong
tài liệu chuyên sâu về inverse spanning tree.

\subsection{Tiểu kết}

Roads là một bài tối ưu hóa mà quy hoạch động hay tham lam không dễ áp dụng
trực tiếp. Sau khi phân tích các tính chất của MST, ta thu được hệ bất đẳng
thức then chốt; từ đó nhận ra cấu trúc giống nhãn khả thi của KM và đưa bài
về ghép cặp tốt nhất. Khi không cần phương án, bài toán còn có thể được nhìn
từ nhiều góc: ghép cặp, dòng chi phí lớn nhất, hoặc ghép có trọng số trên
đỉnh. Việc thay đổi mô hình giúp độ phức tạp giảm dần.

\section{Kết luận}

Qua các ví dụ trên, có thể thấy ứng dụng của ghép cặp đồ thị hai phía trong
thi tin học thường không hiện ra ngay từ đề bài. Ta cần đào sâu bản chất bài
toán, xây dựng đồ thị hai phía thích hợp, rồi dùng ghép cặp để giải. Ngay cả
khi đã nhận ra ghép cặp, cũng không nên chỉ máy móc dùng thuật toán kinh điển;
cần tận dụng tính chất riêng của đề để biến đổi thuật toán và đạt hiệu quả cao
hơn.

Các bài thi tin học nói chung đều đòi hỏi quan sát kỹ đề, dựng mô hình toán
học phù hợp, tiếp tục phân tích mô hình để thấy bản chất, mạnh dạn phỏng đoán
và chuyển hóa mô hình, rồi thiết kế thuật toán thích hợp.

\section*{Lời cảm ơn}

Xin chân thành cảm ơn thầy Xiang Qizhong đã hướng dẫn trong học tập; cảm ơn
Ren Kai, Hu Weidong và Zhou Yuan đã giúp đỡ; cảm ơn Xiao Xiangning và Zhou
Gelin đã góp ý quý báu cho bài viết.

\section*{Tài liệu tham khảo}

\begin{itemize}
  \item Ravindra K. Ahuja và James B. Orlin, \textit{A Faster Algorithm for
  the Inverse Spanning Tree Problem}.
  \item Duan Fanding, tài liệu về thuật toán SPFA nhanh cho đường đi ngắn
  nhất.
\end{itemize}

\appendix
\section{Ghi chú về phát biểu gốc}

Bản gốc kèm phát biểu tiếng Anh của hai bài Saratov State University:
\textit{Railway Communication} và \textit{Roads}. Phần thân bài đã dịch và
tóm lược các yêu cầu cần thiết cho phân tích thuật toán; phụ lục này không
lặp lại toàn bộ phát biểu dài.

\end{document}
